\scanpage{23}
\mspage{10}

\diagram{\begin{tikzcd}[ampersand replacement=\&,column sep=large]
 \&X_{I'}\arrow[r,"\tau_X"]\arrow[d,"\alpha=u_{I'}"']\arrow[dr,phantom,"\text{cocart}" description]
 \&X_{I''}\arrow[d,"\beta"']\arrow[dd,bend left=43,"u_{I''}"]\\
 (*)\&Y_{I'}\arrow[r,"\widetilde\tau_Y"]\arrow[dr,dashed,no head,"\tau_Y"']
 \&\widetilde Y_{I''}\arrow[d,"\widetilde u"'] \&\text{(défini comme \unsure{somme} \unsure{amalgamée})}\\
 \&\&Y_{I''}
\end{tikzcd}}

Alors on a \deleted{$\unsure{\alpha},\beta,\widetilde u$.}

\[
\left\{\begin{aligned}
 &u_{I'},u_{I''}\in C\quad (\text{i.e. pour tt }J\text{ ouvert \unsure{dans|de} }I,\!\\
 &\qquad \inserted{[\text{donc aussi }\beta\in C]}\quad u_J:X_J\mathbin{\unsure{-|\to}}Y_J\text{ est dans }C)^{*}\\
 &\text{et }\deleted{\beta,\unsure{\widetilde u}}\ \tau_X,\tau_Y\in C\quad(\text{i.e. pour }I'\subset I''\\
 &\qquad\text{inclusion d’ouverts, et }X\ C\text{-cofibrant dans}\\
 &\qquad\unsure{\mathcal M|\mathscr M}(I^\circ),\ \unsure{\ldots|\text{on a}}\ X_{I'}\mathbin{-}X_{I''}\text{ dans }C)^{**}\\
 &(\text{\unsure{donc aussi|d’où encore} }\deleted{\unsure{\beta},}\widetilde\tau_Y\in C)\\
 &\deleted{\text{\unsure{et de plus}}}\ \text{et \unsure{enfin}, même}\\
 &\boxed{\widetilde u\in C}\,.
\end{aligned}\right.
\]

% slanted
\margin{$^{*}$c’est bien sûr contenu dans prop.~2, appliqué à $u_\bullet|J$}

% slanted
\margin{C’est pour cette assertion $\tau_X,\deleted{\unsure{\tau_Y}}\in C$ (d’où $\widetilde\tau_Y,\tau_Y\in C$ moyennant $\widetilde u\in C$) qu’on a besoin de l’hypothèse $X_\bullet\in\mathcal M(I^\circ)_C$. Mais $\widetilde u\in C$ vrai \unsure{sans} cette hypothèse}

\deleted{\unsure{L}} \uline{NB} Le cas de la prop.~2 \inserted{sous la forme cas a) (page 9),} \deleted{est \unsure{obtenu}}, (i.e. \unsure{c’est} le fait que \inserted{les} $(u_{I'},u_{I''}\ldots$ sont dans $C$) est obtenu en prenant $I''=I,\ I'=\varnothing$, le diagramme devient

\diagram{\begin{tikzcd}[ampersand replacement=\&,row sep=small]
 \varnothing\arrow[r]\arrow[d,no head,"\mathrm{id}"']\arrow[dr,phantom,"\text{cocart}" description]\& X_I\arrow[d,"\mathrm{id}"]\\
 \unsure{\varnothing}\arrow[r]\&X_I\arrow[d,"\widetilde u="' ,"u_I"]\\
 \&Y_I
\end{tikzcd}}

et la conclusion est donc $u_I\in C$. \inserted{$^{*}$} La \uline{définition} de $C_I$, i.e. la condition $u\in C_I$, est \unsure{retrouvée} en exigeant $\widetilde u\in C$ \inserted{\unsure{dans les cas|du cas}} ci-dessus, dans les cas $I''=I_{/i_0}$, $I'=I_{<i_0}$. Il est remarquable

% margin note written sideways
\margin{$^{*}$ Le cas $(\deleted{\unsure{X}}Y_{I'}\xrightarrow{\tau_{\unsure{Y|X}}}\deleted{\unsure{X}}Y_{I''})\in C$ est obtenu en prenant $X_\bullet=\varnothing_\bullet$, on aura
\diagram{\begin{tikzcd}[ampersand replacement=\&,column sep=tiny,row sep=tiny]
 \varnothing\arrow[r,"\approx","\tau_X"']\arrow[d]\arrow[dr,phantom,"\text{cocart.}" description]\&\varnothing\arrow[d,no head]\\
 Y_{I'}\arrow[r,"\mathrm{id}"]\arrow[dr,"\tau_Y"']\&\widetilde Y_{I''}=Y_{I'}\arrow[d,"\widetilde u={}"' ,"\tau_Y"]\\
 \&Y_{I''}
\end{tikzcd}}}
% archive sheet number: 22

